\documentclass[dvipdfmx]{FITpaper}

\usepackage{booktabs}
\usepackage{flushend}
\usepackage{graphicx}
\usepackage{tabularx}
\usepackage{url}

\newcommand{\doublequotation}[1]{``#1''}
\newcommand{\model}[1]{\mbox{\textrm{#1}}}
\newcommand{\field}[1]{\mbox{\textrm{#1}}}

\makeatletter
\renewcommand{\cite}[1]{[\@nameuse{cite@#1}]}
\expandafter\def\csname cite@clouseau\endcsname{1}
\expandafter\def\csname cite@sysmon\endcsname{2}
\expandafter\def\csname cite@deeplog\endcsname{3}
\expandafter\def\csname cite@react\endcsname{4}
\expandafter\def\csname cite@atlas\endcsname{5}
\expandafter\def\csname cite@atlasv2\endcsname{6}
\makeatother

\begin{document}


% 和文タイトル
\jtitle{SOC調査起点からの証跡付き行動列復元に向けた\\
CLOUSEAU型ログ探索手法の評価}

% 英文タイトル
\etitle{
  Evaluation of CLOUSEAU-based Log Investigation \\
  for Evidence-backed Behavior Reconstruction from SOC Clues
}

% 著者数
\authors{cc}

% 和文著者名
\jauthors{
  著者名 \affmark{1} &
  共著者名 \affmark{2}
}

% 欧文著者名
\eauthors{
  Author Name &
  Coauthor Name
}

% 所属
\afftext{1}{
    所属名, \\
    Affiliation
}
\afftext{2}{
    所属名, \\
    Affiliation
}

\maketitle



\section{はじめに}

Security Operation Center（SOC）におけるインシデント分析では，検知アラートを受け取った後に，端末上で実際に何が起きたかを説明可能な形で再構成する必要がある．
アラート名や短い検知理由は調査の入口として有用である一方，それだけでは親子プロセス，コマンドライン，ファイル・レジストリ・ネットワーク操作，それらを支えるログ行の対応関係を十分に説明できない．
調査者が最終的に必要とするのは，単なる悪性・良性のラベルではなく，観測証跡に基づき，どの主体が，どの対象に，どの順序で操作したかを示す行動列である．

本研究が対象とするのは，マルウェア検知そのものではなく，検知後のエンドポイント調査である．
特に，SOCが少数の手がかりから調査を開始する状況を想定する．
入力は，EDRのalert要約，またはhost，process，time window程度のprocess/time clueであり，出力は証跡付きの行動ステップ列である．
この設定では，CLOUSEAUのようなLLM階層型調査エージェント\cite{clouseau}が，自然言語の手がかりから探索方針を作り，SQLite化されたログに対して反復的に問い合わせを行うことで，関連する観測行を集められる可能性がある．
しかし，アラート理由を言い換えただけの出力，近傍の無関係な行動の混入，証跡で支えられない推定，行動順序の誤りが生じる危険もある．

本稿では，ATLASv2 benign H1 / benign-1由来のWindowsホストログを用い，23個の良性セキュリティ行動chainを対象に，CLOUSEAUを基盤とするエージェント型ログ探索手法を評価する．
ここで既存研究との立ち位置を明確にしておく．
CLOUSEAUは，Point-of-Interest（POI）から攻撃調査を開始し，Chief Inspector，Investigator，Question-Answering agentが協調してattack narrativeを復元する階層型multi-agent frameworkである\cite{clouseau}．
ATLASはaudit logから攻撃ストーリを復元する攻撃調査支援手法であり\cite{atlas}，ATLASv2はその攻撃エンゲージメントと通常業務活動を拡張したデータセットである\cite{atlasv2}．
本研究はCLOUSEAUやATLASを置き換える攻撃調査アルゴリズムではない．
CLOUSEAUの調査パイプラインとATLASv2のログ環境を用い，SOCの検知後調査を模した行動列復元ベンチマークとして，アラート要約への依存度，低レベルtelemetryからの復元可能性，証跡同定の成否を測る研究である．

本稿の貢献は次の3点である．
第一に，CLOUSEAUのPOI起点調査をSOCの検知後行動復元へ読み替え，行動item，critical evidence，sequence order，candidate precision，overclaimに分解した評価枠組みとして定義する．
第二に，MITRE ATT\&CKとLOLBAを参考に実務上調査対象になり得るprocessを選び，CBC alertで偽陽性として現れた時刻窓から23 usecaseを作る．
これらを3段階の入力条件で実行し，アラート支援条件，process/time clue条件，アラート要約行を隠したtelemetry中心条件を比較する．
第三に，ATLASv2由来の単一Windows良性環境における検知後分析として射程を限定し，明示的実行チェーン，複数段ツールチェーン，意味解釈型チェーンごとの失敗傾向を分析する．



\section{関連研究}

\subsection{検知後調査としての行動列復元}

エンドポイント調査では，Windows Securityログ，Sysmon，EDR telemetry，DNS，ブラウザ履歴など，粒度や記録目的の異なるログを横断する．
Sysmonはプロセス生成，親プロセス，コマンドライン，ネットワーク接続，レジストリ操作などを記録できる\cite{sysmon}．
EDRは検知理由やプロセス情報をまとめたalert summaryを提供する場合がある．
これらは調査に有用だが，alert summaryは観測された全行動を代表するものではない．
したがって，最終出力はアラート文面ではなく，ログ上で確認できる行動主張に限定する必要がある．

本研究の課題は，異常検知モデルの精度評価ではない．
DeepLogのような系列異常検知研究は，ログ系列から異常を検出・診断することに焦点を置く\cite{deeplog}．
一方，本研究は検知後のSOC調査において，調査者が検証可能な証跡付き行動列を復元できるかを問う．
この違いは評価単位にも表れる．
本研究では自然文の類似度ではなく，主体，操作，対象，コマンドライン，証跡行などの必須itemを採点する．

\subsection{LLMエージェントによるログ探索}

LLMエージェントは，初期手がかりから仮説を立て，外部環境への問い合わせ結果に応じて次の探索を変更できる．
ReActは推論と行動を交互に進める枠組みを示しており\cite{react}，ログ調査でも，時刻範囲，プロセス木，ファイル・レジストリ対象を段階的に確認する流れと親和性がある．

セキュリティ分野では，CLOUSEAUが階層型multi-agentによる自律的な攻撃調査を提案している\cite{clouseau}．
CLOUSEAUは，raw incident logsをLLMが扱いやすい構造化形式へ変換し，QA agentを通じてログを自然言語的に検索可能にする．
さらに，Investigator agentとChief Inspector agentが調査leadの生成，証跡収集，attack storyの統合を行う．
本研究はこの設計思想に基づくが，評価対象を「攻撃story全体の復元」から「SOCアラートやprocess/time clue周辺の良性を含む行動chainの証跡付き復元」へ移す．
そのため，最終判定の正否ではなく，観測ログで支えられた行動itemと過剰主張を採点する．

\subsection{ATLAS/ATLASv2の位置づけ}

ATLASは，audit logに基づいて攻撃ストーリを復元するため，因果依存グラフと系列学習を組み合わせた攻撃調査支援手法である\cite{atlas}．
ATLASv2は，ATLASの攻撃シナリオを再現しつつ，通常業務活動とSysmonおよびCarbon Black Cloud由来のログを拡充したデータセットとして報告されている\cite{atlasv2}．
本研究では，ATLASv2を「攻撃検知済みの答えを読むデータ」ではなく，「SOCでアラートやtelemetryを受け取った後，端末上の行動を証跡付きで説明する能力を測るデータ」として利用する．
この立ち位置により，本研究の対象インシデント分析は，侵害全体の帰属や最終判定ではなく，アラート周辺の端末行動復元であると明確化される．



\section{提案手法}

\subsection{パイプライン}

図\ref{figure:pipeline}に，本研究で用いるCLOUSEAU型のエージェント型ログ探索パイプラインを示す．
入力は，CBC alert rows，またはhost/process/time windowである．
CLOUSEAU本来のPOI起点調査\cite{clouseau}にならい，Chief agentは初期手がかりから調査方針を生成する．
Investigation agentは，確認すべき事実をQAAgent / SQL expertへの質問に変換する．
QAAgent / SQL expertは，SQLite化されたエンドポイントログに問い合わせ，timestamp，source stream，PID/PPID，process tree，command line，registry path，source row idなどを返す．
Final synthesisは観測結果を統合し，主行動列，補助証跡，limitationsを出力する．

\begin{figure}[ht]
    \centering
    \fbox{
    \begin{minipage}{0.92\linewidth}
        \small
        \centering
        SOC clue / CBC alert / process-time clue\\
        $\downarrow$\\
        Chief agent: 調査方針を生成\\
        $\downarrow$\\
        Investigation agent: 確認事項を質問へ変換\\
        $\downarrow$\\
        QAAgent / SQL expert: 観測行を取得\\
        $\downarrow$\\
        Final synthesis: 主行動列と近傍行動を分離\\
        $\downarrow$\\
        行動ステップ列 + evidence + limitations
    \end{minipage}}
    \caption{SOC調査起点から証跡付き行動列を復元するパイプライン}
    \label{figure:pipeline}
\end{figure}

本手法の設計上の要点は，仮説生成と観測ログ検索を分離することである．
Chief agentは直接SQLを書かず，Investigation agentが観測済みの値に基づいて次の確認事項を作る．
DBに問い合わせるのはQAAgent / SQL expertであり，Final synthesisは観測行で支えられる行動だけを主系列に含める．
これにより，alert nameやreasonの言い換えを行動復元と誤認しないようにする．

\subsection{出力と評価単位}

出力は，\field{steps}，\field{sequence}，\field{evidence}，\field{nearby}，\field{limitations}からなる．
\field{steps}は主体，操作，対象，コマンドライン，証跡を持つ主行動ステップである．
\field{sequence}は観測された操作を短く並べた系列であり，アラート名ではなくログ上の操作を表す．
\field{evidence}はsource stream，row id，timestamp，PID，command lineなどの根拠である．
\field{nearby}は時刻や対象が近いが主系列に含めない行動を明示する．

採点では，gold chainを行動itemに分解する．
必須itemは，ログから直接確認できる主体，操作，対象，コマンドライン，critical evidenceである．
観測行だけから主張できない悪性判定，ユーザ意図，業務目的，ファイル内容の推定は必須itemに含めない．
候補出力の主張は正規化され，gold itemと照合される．
主指標は，行動item recall，critical evidence recall，sequence order，candidate claim precision，overclaim slot countである．



\section{実験設定}

\subsection{データとchain}

実験には，ATLASv2 benign H1 / benign-1環境をCLOUSEAUのログ探索形式へ変換したWindowsホストログを用いた．
評価単位は23個のbehavior chainである．
usecaseは，単にログ中の行を機械的に切り出したものではなく，実務上SOCで調査対象になり得るprocessを起点に選定した．
まず，MITRE ATT\&CKやLOLBAを参考に，攻撃でも正常運用でも使われ得るWindows上のprocessや一般ツールを候補化した．
次に，ATLASv2内のCBC alertで実際に良性環境中に検知された時刻窓を抽出し，CBC alert，focus process，周辺telemetryから一つの行動背景を説明できるものをchain化した．
この手順により，CBC alertを起点にSOC調査が始まるが，最終的には良性行動として説明すべき23 usecaseを得た．

chainは，DNS packet capture batch，Python SimpleHTTPServer network，Sublime Python script execution，cmd.exe周辺操作，Discord Run Key registryの5種類のraw chain typeから構成される．
分析時には，これらを明示的実行チェーン，複数段ツールチェーン，意味解釈型チェーンの3群にまとめる．
表\ref{table:usecases}に，usecaseの内訳を示す．

\begin{table}[ht]
    \centering
    \caption{usecase選定結果}
    \label{table:usecases}
    \begin{tabular}{lrr}
        \toprule
        行動背景 & Chains & Group \\
        \midrule
        Python簡易Webサーバ & 13 & 明示的実行 \\
        DNS/通信ログ収集batch & 5 & 複数段ツール \\
        Sublime経由Python実行 & 2 & 複数段ツール \\
        cmd.exe周辺操作 & 2 & 明示的実行 \\
        Discord Run key登録 & 1 & 意味解釈 \\
        \bottomrule
    \end{tabular}
\end{table}

明示的実行チェーンは，親子プロセスとコマンドラインが比較的直接的に観測できるchainである．
複数段ツールチェーンは，batch，packet capture tool，editor，script，file/network side effectなど，複数の道具や中間ステップを結ぶ必要があるchainである．
意味解釈型チェーンは，Discord/Update.exeとRun Keyのように，アプリケーション固有またはレジストリ固有の意味解釈を必要とするchainである．

\subsection{3段階の入力条件}

表\ref{table:stages}に，3つのstageを示す．
Stage 1はCBC alert rowsを入力に含めるアラート支援条件である．
Stage 2はhost，focus process，time windowのみを入力とし，DB内にはalert summary rowsを残す．
Stage 3はStage 2と同じ初期入力だが，検索対象からCBC alert summary rowsを除外する．
ただし，CBC EDR/NGAV event telemetry，Security，Sysmon，DNS，browser historyは残す．

\begin{table}[ht]
    \centering
    \caption{入力条件}
    \label{table:stages}
    \begin{tabularx}{\linewidth}{lXX}
        \toprule
        Stage & 入力 & 意味 \\
        \midrule
        Stage 1 & CBC alert rows & アラート支援の復元 \\
        Stage 2 & host/process/time & DB内の要約行は利用可能 \\
        Stage 3 & host/process/time & alert summary rowsを検索対象から除外 \\
        \bottomrule
    \end{tabularx}
\end{table}

この設計により，初期入力の強さと，DB内の高密度なアラート要約行の有無を分けて観察できる．
Stage 2とStage 3の差は，alert summaryが最終出力として使われるかではなく，探索経路や証跡同定の足場として機能するかを調べるためのものである．

\subsection{実行と採点}

表\ref{table:scope}に，本稿で用いる評価系列を整理する．
正式な性能比較に用いるのは，23 chain，3 stage，3 setからなる各モデル207 runである．
3 setのうち2つは23 chainとして実行したreplicateであり，残る1つは旧27 chain実行から同じ23 chain IDだけを抽出したsetである．
\model{gpt-5.5}は23 chain，3 stageの69 runを実行したが，構造化JSON契約を満たさないraw text出力を後処理でsalvageしたため，正式比較とは分けて扱う．

\begin{table}[ht]
    \centering
    \caption{評価系列と母数}
    \label{table:scope}
    \begin{tabular}{lrrrr}
        \toprule
        Use & Chains & Stages & Sets & Runs \\
        \midrule
        正式比較 & 23 & 3 & 3 & 207/model \\
        \model{5.5} raw & 23 & 3 & 1 & 69 \\
        \bottomrule
    \end{tabular}
\end{table}

正式比較の対象は\model{gpt-4.1-mini}と\model{gpt-5.4-mini}である．
各モデルについて，23 chainを3 stageで3 set評価し，207 runを得た．
2モデル合計で414 runを採点し，\model{gpt-5.5}は69 runのraw text salvageとして別枠で確認した．
採点対象のgoldは，起点となる観測行の特定，process actionまたはobject operationへの分解，必須item定義，critical evidence slot固定，gold chainに属さない近傍行動の除外という手順で作成した．

採点では，単に似た説明文が出ているかではなく，行動を構成する最小単位がログ証跡に接続しているかを重視した．
たとえば，``Pythonが実行された''という説明だけでは，主体，実行対象，親子関係，コマンドライン，時刻のいずれを根拠としているかが曖昧である．
そこで各stepを，主体，操作，対象のaction itemに分け，それとは独立に，該当stepを支えるcritical evidence itemを定義した．
この分離により，モデルが行動の概略を当てたが証跡を示せない場合と，証跡行を拾ったが行動列として整理できない場合を区別できる．
sequence orderは，隣接するgold stepの前後関係が候補出力で保たれた割合として数えた．

また，候補出力の品質を見るため，candidate claim precisionとoverclaimも記録した．
candidate claim precisionは，モデルが主張した候補claimのうちgold chainに対応するものの割合である．
overclaimは，近傍イベント，時刻窓内に存在するが主系列ではない操作，または証跡からは断定できない因果関係を，chainの一部として述べた件数である．
SOC調査支援では広めの候補提示が有用な場面もあるため，overclaimをただちに失敗とは見なさない．
しかし，本研究で評価したいのは「証跡に基づく行動列復元」であるため，正解itemを多く拾うrecallだけでなく，主系列と補助情報の境界をどれだけ保てるかを併せて評価した．

本稿の23 chain集計では，各stageあたり65 behavior stepsを基準とし，1 setでは3 stage合計195 step，3 set合計では585 stepを評価する．
各stepは主体，操作，対象の3つのaction itemと，1つのcritical evidence itemを持つ．
したがって，正式比較の分母は各モデルについてaction itemが1755，critical evidenceが585，orderが378である．
なお，第3 setは旧27 chain実行から現在の23 chainに対応する行だけを抽出したものであり，完全な新規replicateではない．
この点は，3 set平均の分散解釈に制約を与える．

採点は単一の自動judgeだけで確定しない．
候補出力とgoldを並べたレビュー行を作成し，2回の独立レビューを行い，一致した項目を採用した．
判断が割れた項目はconflict queueへ入れ，追加レビューまたは保守的な採択で確定した．
表の値は，表現やJSON fieldの完全一致ではなく，必要な内容が候補出力に含まれるかを見るcontent inclusion方針で集計した．

\subsection{評価範囲と除外条件}

本研究の評価範囲は，alertまたはprocess/time clueを起点として，端末上の具体的な行動chainを復元することである．
マルウェア分類，攻撃キャンペーン attribution，MITRE ATT\&CK technique labelの推定，検知ルールの良否評価は主目的ではない．
この限定は，ATLASv2とCLOUSEAUを利用する本研究の立ち位置を明確にするためである．
ATLASv2は攻撃調査に近いログ環境を提供し，CLOUSEAUはPOIから証跡を探索してnarrativeを組み立てる枠組みを与える．
本稿はその上で，SOCがアラート後に行う「何が起きたかを証跡付きで説明する」作業に対象を絞る．

除外したものも明示しておく．
第一に，出力が最終的に悪性か良性かを判定できたかは採点しない．
対象chainには良性操作が含まれるため，悪性判定を正解とする評価は本研究の目的とずれる．
第二に，長い自然文として読みやすいかは補助的に確認するが，主指標にはしない．
読みやすい報告書であっても，証跡IDやイベント行に接続できなければSOCの検証対象としては不十分である．
第三に，Stage 3はalert summary rowsを隠す条件であり，すべてのEDR telemetryを取り除く条件ではない．
このため，Stage 3の結果は，低レベルログだけで常に十分であることではなく，高密度なsummaryを欠いても残存telemetryから復元可能な部分があることを示す．

公開用パッケージには，case metadata，chain summary，gold index，runごとのモデル出力JSON，採点rubric，集計CSV/JSONを含めた．
一方，元の\texttt{incident.db}，\texttt{scenario.db}，EVTX，外部データセット本体は含めない．
したがって，本稿の表とrun-level挙動は検査可能であるが，完全な再実行にはローカルのATLASv2/CLOUSEAUデータベース環境が必要である．



\section{結果と考察}

\subsection{全体結果}

表\ref{table:overall}に，全体結果を示す．
\model{gpt-5.4-mini}は207 run全体で行動item recall 79.8\%（1400/1755），critical evidence recall 70.3\%（411/585）を示した．
一方，\model{gpt-4.1-mini}は行動item recall 46.6\%（818/1755），critical evidence recall 19.5\%（114/585）であった．
この比較は23 chainの3 set集計であり，第3 setに旧27 chain実行からの23 chain抽出を含む点に注意が必要である．

\begin{table}[ht]
    \centering
    \caption{全体結果}
    \label{table:overall}
    \begin{tabular}{lrrr}
        \toprule
        Model & Action & Crit. ev. & Order \\
        \midrule
        \model{gpt-4.1-mini} & 46.6\% & 19.5\% & 20.1\% \\
        \model{gpt-5.4-mini} & 79.8\% & 70.3\% & 55.3\% \\
        \bottomrule
    \end{tabular}
\end{table}

注目すべき点は，\model{gpt-5.4-mini}でもcandidate claim precisionは58.4\%に留まったことである．
これは，強いモデルが正解itemを多く拾える一方で，候補出力に近傍行動や重複主張も残すことを意味する．
SOC支援としては広めに候補を出す価値があるが，論文上の行動列復元としては，主系列と補助証跡の分離が依然として課題である．

表\ref{table:claimquality}に，claim qualityの補助指標を示す．
\model{gpt-4.1-mini}はcandidate precisionが36.6\%で，overclaimが1190件であった．
一方，\model{gpt-5.4-mini}はcandidate precisionを58.4\%まで改善し，overclaimも651件まで減らした．
ただし，半数近い候補claimがgold chainに直接対応しない点は残っている．
この結果は，LLMが調査文脈で有用な候補生成器になり得る一方，そのまま最終報告書として使うには，claimの昇格・棄却を制御する仕組みが必要であることを示す．

\begin{table}[ht]
    \centering
    \caption{候補claim品質}
    \label{table:claimquality}
    \begin{tabular}{lrr}
        \toprule
        Model & Precision & Overclaim \\
        \midrule
        \model{gpt-4.1-mini} & 36.6\% & 1190 \\
        \model{gpt-5.4-mini} & 58.4\% & 651 \\
        \bottomrule
    \end{tabular}
\end{table}

この差は，単なる表現力の差だけでは説明しにくい．
\model{gpt-4.1-mini}は，時刻窓内に存在するイベントを広く拾い，それらをchainの一部として結合しがちであった．
これに対して\model{gpt-5.4-mini}は，親子プロセス，ファイルパス，コマンドライン，引数，通信先などの証跡を比較して，主系列らしい候補を選ぶ傾向が強かった．
しかし，Stage 3や複数段ツールチェーンでは，packet capture toolの実行，script fileの編集，HTTP server起動，browser accessなどが時刻的に近接するため，主系列と周辺操作を完全に分けることは難しい．
この点が，recallの高さとprecisionの伸び悩みが同時に現れる理由である．

\subsection{Stage別結果}

表\ref{table:bystage}にStage別結果を示す．
\model{gpt-5.4-mini}はStage 1で79.1\%（463/585），Stage 2で80.0\%（468/585），Stage 3で80.2\%（469/585）の行動item recallを示した．
23 chain最終集計では各stageの分母を揃えており，Stage 3だけ小さい分母で見せてはいない．
したがって，alert summary rowsを隠した条件でも，残存telemetryから同程度のaction itemを復元できたと読める．

\begin{table}[ht]
    \centering
    \caption{Stage別結果（23 chain，3 set）}
    \label{table:bystage}
    \begin{tabular}{llrr}
        \toprule
        Model & Stage & Action & Crit. ev. \\
        \midrule
        \model{4.1-mini} & S1 & 63.9\% & 29.2\% \\
        \model{4.1-mini} & S2 & 41.4\% & 12.8\% \\
        \model{4.1-mini} & S3 & 34.5\% & 16.4\% \\
        \model{5.4-mini} & S1 & 79.1\% & 57.4\% \\
        \model{5.4-mini} & S2 & 80.0\% & 75.4\% \\
        \model{5.4-mini} & S3 & 80.2\% & 77.9\% \\
        \bottomrule
    \end{tabular}
\end{table}

一方で，\model{gpt-5.4-mini}のStage 2/3のorderは54.0\%であり，Stage 1の57.9\%よりやや低い．
process/time clueだけで探索を開始する場合，必要なitemに到達できても，それらを正しい順序または因果的構造として並べることは難しい．
Stage 3ではclaim precisionが59.4\%で維持された一方，overclaimは223件であり，候補主張の量は依然として多い．
alert summaryがない条件では，低レベルtelemetryを広く探索するため，主系列に含めるべき行動と近傍行動の境界が曖昧になりやすい．

表\ref{table:stageorder}に，\model{gpt-5.4-mini}のStage別orderとclaim precisionを示す．
Stage 1ではalert rowsが起点と順序の足場を与えるため，orderは57.9\%であった．
Stage 2では，DB内にalert summary rowsが残っているにもかかわらず，入力がhost/process/time clueへ弱まることでorderが54.0\%となった．
つまり，探索対象にsummaryが存在することと，最初からalert文脈が入力されることは等価ではない．
Stage 3でもorderは54.0\%であり，precisionは59.4\%であった．
これは，summaryを検索できないことで，モデルが低レベル証跡に寄った順序付けを試みる一方，候補claimを多めに保持するためと考えられる．

\begin{table}[ht]
    \centering
    \caption{\model{gpt-5.4-mini}のStage別補助指標}
    \label{table:stageorder}
    \begin{tabular}{lrrr}
        \toprule
        Stage & Order & Precision & Overclaim \\
        \midrule
        S1 & 57.9\% & 55.9\% & 215 \\
        S2 & 54.0\% & 59.7\% & 213 \\
        S3 & 54.0\% & 59.4\% & 223 \\
        \bottomrule
    \end{tabular}
\end{table}

この結果から，alert summary rowsには少なくとも二つの働きがあると解釈できる．
第一に，調査開始時の意味的な手掛かりとして，どのprocessやobjectを中心に見るべきかを制限する働きである．
第二に，探索途中で発見した低レベルイベントを，chainのどの位置に置くかを決める境界条件としての働きである．
Stage 2は後者を利用できるが前者が弱く，Stage 3は両者のうちsummary由来の足場をさらに失う．
そのため，Stage 3でrecallが高くても，それはsummary不要という意味ではない．
むしろ，summaryがない場合には，候補claimの剪定と順序制約がより重要になることを示している．

\subsection{Scenario別結果}

表\ref{table:scenario}に\model{gpt-5.4-mini}のscenario group別結果を示す．
明示的実行チェーンではaction recall 86.5\%，critical evidence recall 82.4\%，order 76.4\%であり，親子プロセスやコマンドラインが直接観測できるchainでは高い性能を示した．
複数段ツールチェーンではaction recall 74.6\%を示したが，orderは43.1\%，claim precisionは46.9\%に低下した．
意味解釈型チェーンは1 chainを3 stage，3 setで評価した9 runのみであり，数値は傾向の示唆に留まる．
その範囲では，critical evidence recall 48.1\%，order 33.3\%となり，registry/app固有の意味解釈が難所である可能性を示した．

\begin{table}[ht]
    \centering
    \caption{\model{gpt-5.4-mini}のscenario別結果}
    \label{table:scenario}
    \begin{tabular}{lrrr}
        \toprule
        Group & Action & Crit. ev. & Order \\
        \midrule
        明示的実行 & 86.5\% & 82.4\% & 76.4\% \\
        複数段ツール & 74.6\% & 60.2\% & 43.1\% \\
        意味解釈 & 64.2\% & 48.1\% & 33.3\% \\
        \bottomrule
    \end{tabular}
\end{table}

\subsection{gpt-5.5と運用コスト}

\model{gpt-5.5}については，正式なJSON契約を満たした207 runの3 set採点結果がないため，表\ref{table:overall}の正式比較には含めない．
23 chainに対するraw text salvageでは高い復元傾向が見られたが，raw textから後処理でclaimを抽出した単一replicateであり，正式JSON評価とは同列に扱えない．
したがって，本稿では\model{gpt-5.5}を性能結論ではなく，出力契約と費用制御を改善した次実験の候補として位置づける．

表\ref{table:cost}に，運用コストと実行時間の補助集計を示す．
\model{gpt-5.5 low raw}は，69 runで入力1.83M tokens，出力0.73M tokensを消費し，推定費用は31.00 USDであった．
平均実行時間は9.32分で，直列合計は10.72時間，4並列相当では2.68時間である．
\model{gpt-5.4-mini}の補助ログは207 runで4.35 USD，平均2.05分であり，費用・時間の両面で軽い．
したがって，\model{gpt-5.5}は復元品質の上限を探る参照点として有用だが，出力契約失敗と高い出力token費用を考えると，全量実験の標準モデルとしては扱いにくい．

\begin{table}[ht]
    \centering
    \caption{モデル別コスト・token・時間}
    \label{table:cost}
    \begin{tabular}{lrrrr}
        \toprule
        Model & Runs & Cost & Tok. & Avg min \\
        \midrule
        \model{4.1-mini} & 207 & \$3.64 & 7.40M & 6.81 \\
        \model{5.4-mini} & 207 & \$4.35 & 2.67M & 2.05 \\
        \model{5.5 raw} & 69 & \$31.00 & 2.56M & 9.32 \\
        \bottomrule
    \end{tabular}
\end{table}

今後\model{gpt-5.5}を正式採用するには，少数pilotでJSON契約，max output tokens，累積費用停止条件を確認してから，23 chain $\times$ 3 stage $\times$ 3 setの同一rubricで再実行する必要がある．

\subsection{考察}

本研究の結果は，CLOUSEAUとATLASv2を組み合わせることで，検知後の端末行動復元を評価できることを示している．
CLOUSEAUは攻撃narrative復元のための階層型調査エージェントであり，ATLAS/ATLASv2は攻撃調査文脈のデータ・手法として位置づけられてきた．
本研究では，CLOUSEAUのPOI起点探索をSOCのアラート後分析に引き寄せ，ATLASv2のログ環境上で，検知ラベルの分類ではなく，アラートやprocess/time clueからログ証跡で支えられた行動chainを構成する能力を測った．

良性chainを対象にした点も重要である．
SOCでは，悪性攻撃だけでなく，管理操作，開発作業，ツール実行，アプリケーション更新なども調査対象になる．
これらはアラートを発生させたり，攻撃手法に似たtelemetryを持ったりする．
したがって，良性行動を証跡付きで説明できることは，誤検知低減や調査記録作成に直結する．

Stage 1はアラート情報を入力に含み，Stage 2は入力から外すがDBには残し，Stage 3は検索対象からalert summary rowsを隠す．
Stage 3で\model{gpt-5.4-mini}が分母を揃えた条件でも高いrecallを示したことは，出力がalert summaryの単純な言い換えではなく，残存telemetryからも多くの行動itemを復元できたことを示す．
ただし，Stage 3でもoverclaimは残るため，alert summaryは探索の近道や境界設定として機能していた可能性がある．

\subsection{CLOUSEAUとの差分}

本研究はCLOUSEAUに強く依存しているため，CLOUSEAUとの違いを曖昧にすると，貢献が見えにくくなる．
CLOUSEAUは，POIを起点としてChief Inspector，Investigator，QA agentが協調し，攻撃narrativeを再構成する階層型multi-agent frameworkである．
本稿はこの枠組みを，攻撃キャンペーン全体のnarrative復元ではなく，SOCのアラート後調査で頻出する短いbehavior chain復元に適用し，入力条件と評価粒度を変えて検証した．
特に，良性chainを含む点，Stage 1--3でalert summaryの有無を分ける点，action itemとcritical evidence itemを分離して採点する点が異なる．

この違いは，研究目的の違いでもある．
CLOUSEAUの主眼は，攻撃を説明するnarrativeをエージェントが構築できるかである．
本研究の主眼は，SOC調査者が最終判断を行う前段として，モデルがログ上の行動単位と根拠行をどれだけ正確に揃えられるかである．
そのため，本稿では自然文としてもっとも説得的なstoryよりも，gold step，critical evidence，order，overclaimを重視した．
これは，実運用のインシデント分析では，華麗な説明よりも，後から検証可能な証跡列が必要になるためである．

\subsection{失敗傾向}

失敗は三つに整理できる．
第一は，近傍イベントの巻き込みである．
同じ時刻窓に複数の管理操作や開発作業が存在すると，モデルはそれらを一つのchainに結合しやすい．
これはStage 3で特に目立ち，alert summaryによる境界が消えると，関連しそうな低レベルイベントを広く残す傾向が強まる．

第二は，意味解釈型chainでの根拠不足である．
Discord/Update.exeとRun Key registryのようなchainでは，プロセス名やレジストリパスの意味を理解する必要がある．
単にイベント行を拾うだけでは，永続化，更新処理，アプリ固有挙動のどれとして読むべきかが決まらない．
このため，critical evidence recallとorderがともに低下した．

第三は，順序と因果の混同である．
ログ上の時刻順は，行動の因果順と一致するとは限らない．
特に，ファイル作成，プロセス起動，ネットワーク接続，DNS queryが短時間に発生するchainでは，モデルが観測順をそのまま因果順として述べる場合があった．
この問題を解くには，単なる検索結果の列挙ではなく，親子関係，object identity，command line argument，network endpointを束ねる制約付きのsynthesisが必要である．

\subsection{実用上の含意}

実用上，本結果はLLMを単独の自動判定器として使うよりも，調査者が確認すべき候補chainを作る補助器として使う方が現実的であることを示す．
\model{gpt-5.4-mini}は，多くの正解itemとcritical evidenceを拾えるため，初動調査の探索範囲を狭める用途に向く．
一方で，candidate precisionが58.4\%である以上，出力をそのままインシデント報告書へ転記する設計は危険である．
モデル出力は，証跡行，主系列claim，補助claim，棄却候補に分け，調査者が昇格判断を行うUIやworkflowと組み合わせる必要がある．

また，Stage 3の結果は，アラートsummaryが失われた環境でも一定の復元が可能であることを示すが，summaryを不要化できることを意味しない．
むしろ，summaryは調査対象を狭め，overclaimを抑え，順序を安定させる補助情報として有用である．
SOCにおける望ましい運用は，alert summaryを根拠として鵜呑みにすることではなく，summaryを起点に低レベルtelemetryへ降り，モデルに再構成させたchainを再び証跡で検証する循環である．



\section{おわりに}

本稿では，CLOUSEAUを基盤とするエージェント型ログ探索手法により，SOC調査起点からWindowsエンドポイントログを探索し，証跡付き行動列を復元する課題を，ATLASv2 benign H1 / benign-1由来の23 behavior chainで評価した．
\model{gpt-5.4-mini}は207 run全体で行動item recall 79.8\%，critical evidence recall 70.3\%を示した．
alert summary rowsを隠したStage 3でも，分母を揃えた条件で80.2\%の行動item recallを示した．
一方で，candidate claim precision，overclaim，sequence orderには課題が残った．

本研究の立ち位置は，CLOUSEAUの攻撃narrative復元をそのまま再評価するものでも，ATLAS/ATLASv2を用いた攻撃検知評価でもなく，検知後のSOCインシデント分析における行動列復元評価である．
今後は，SQL探索履歴の分析，主系列昇格条件の強化，順序制約を持つsynthesis，意味解釈型chainの拡充により，調査者がそのまま検証できる行動storyの品質を高める．

本評価の限界として，対象は単一のATLASv2 benign環境から構成した23 chainであり，組織横断のSOCログや攻撃キャンペーン全体を代表するものではない．
また，Stage 3はalert summary rowsのみを検索対象から除外する条件であり，EDR telemetry全体を除外した条件ではない．
このため，本稿の結論は「alert summaryなしでも一定の行動復元が可能」という範囲に留め，「任意の低レベルログだけで十分」とは主張しない．

妥当性への脅威は五点ある．
第一に，gold chainの作成はログ構造と実験目的を理解した上で行っており，完全に自動的な正解ではない．
このため，step境界の切り方やcritical evidence slotの選び方に，作成者の判断が入る．
本稿では，主体，操作，対象，根拠行を分けて定義し，23 chain最終集計ではstage間の分母を揃えることで，恣意的な高得点化を避けた．
第二に，評価対象のchainは良性環境に由来するため，攻撃者が意図的に痕跡を隠すケース，複数ホストへ横展開するケース，長時間に分散するキャンペーンは含まれていない．
したがって，本結果は攻撃調査全体の代替ではなく，単一ホストのアラート後分析に近い短いchain復元の結果として読む必要がある．
第三に，本稿には単発LLMや単純時刻窓検索とのベースラインがなく，CLOUSEAU型パイプライン自体の上乗せ効果は分離できていない．
第四に，API version，temperature，top-p，seed，プロンプト全文などの再現情報は公開パッケージ側に依存しており，本文中には十分に収められていない．
第五に，本稿の比較は二つのモデルに限られる．
特に，\model{gpt-5.5}についてはraw text salvageの補助結果とコスト・token・時間の記録を示したが，同一rubricで採点済みの23 chain $\times$ 3 stage $\times$ 3 set正式JSON結果ではない．
未採点または分母の異なる結果を正式比較へ混ぜると，モデル比較として不正確になるためである．
ただし，\model{gpt-5.5}が高い復元率を示した点は，出力契約と費用制御を改善した次実験の有力な候補であることを示す．

今後の実装では，モデル出力を一つの文章として受け取るのではなく，claim graphとして扱う．
各claimに，対応するevent row，親子process関係，object identity，時刻範囲，confidence，gold chain候補への対応を付与し，主系列claimと補助claimを明示的に分ける．
その上で，調査者がclaimを採用，保留，棄却できるUIを用意すれば，高recallな探索と人間によるprecision改善を両立できる．
また，CLOUSEAUのQA agentに相当する検証段を，単なる自然文批評ではなく，未接続claim，証跡不足claim，順序矛盾claimを検出する形式へ拡張することが有効である．
これにより，Stage 3で顕在化したoverclaimを，報告書生成前に抑制できる可能性がある．



\section*{参考文献}

\begin{enumerate}
    \item A. Aldaihan, F. Alotaibi, and S. Maffeis,
    ``Clouseau: A Hierarchical Multi-Agent Approach For Autonomous Attack Investigation,''
    2025. \url{https://www.doc.ic.ac.uk/~maffeis/papers/acsac25.pdf}

    \item Microsoft,
    ``Sysmon - Sysinternals,''
    \url{https://learn.microsoft.com/en-us/sysinternals/downloads/sysmon}
    （参照 2026-06-09）．

    \item M. Du, F. Li, G. Zheng, and V. Srikumar,
    ``DeepLog: Anomaly Detection and Diagnosis from System Logs through Deep Learning,''
    Proc. ACM CCS, pp. 1285--1298, 2017.

    \item S. Yao, J. Zhao, D. Yu, N. Du, I. Shafran, K. Narasimhan, and Y. Cao,
    ``ReAct: Synergizing Reasoning and Acting in Language Models,''
    arXiv:2210.03629, 2022.

    \item A. Alsaheel, Y. Nan, S. Ma, L. Yu, G. Walkup, Z. B. Celik, X. Zhang, and D. Xu,
    ``ATLAS: A Sequence-based Learning Approach for Attack Investigation,''
    Proc. USENIX Security, pp. 3005--3022, 2021.

    \item A. Riddle, K. Westfall, and A. Bates,
    ``ATLASv2: ATLAS Attack Engagements, Version 2,''
    arXiv:2401.01341, 2023.
\end{enumerate}


\end{document}
